\documentclass[instructions.tex]{subfiles}
\begin{document}
 
\chapter{Il faut réparer le tort qu’on a fait, souffrir celui qu’on nous cause, et comment.}
 
\primo On s’aveugle quand on prend le bien d’autrui; mais on s’aveugle bien davantage quand on l’a pris. Je ne suis pas obligé à rendre, dit-on; je me suis consulté; j’ai fait une compensation; d’ailleurs je ne le puis.
 
Si un pasteur avertit un homme de mauvaise foi qu’il est obligé à restituer; s’il s’oppose aux injustices, aux abus d’une paroisse, on ne l’écoute plus, on le méprise, on le persécute. Une décision indiscrète d’un homme sans expérience, un avis artificieusement extorqué d’un avocat qu’on surprend, l’emporteront sur ce qu’on doit aux prêtres du Seigneur, à qui il appartient d’expliquer la loi de Dieu\footnote{Labia enim sacerdotis custodiunt scientiam, et legem requirent ex ore ejus. Malac. 2.}.
 
Les avares et les larrons sont presque tous incorrigibles. Jésus-Christ d’un seul regard toucha le cœur de saint Pierre; Magdeleine entendit le Sauveur, et aussitôt elle fondit en larmes; d’une parole il changea Matthieu; mais ni Judas, ni aucun pharisien ne se convertirent, parce qu’ils étaient avares et larrons.
 
Quand vous verseriez autant de larmes que tous les pénitens, jamais vos injustices ne seront pardonnées, dit saint Augustin, si vous ne restituez\footnote{Non, dimittitur peccatum nisi restituatur ablatum. S. Aug.}. Si vous pouvez restituer, vous y êtes obligé sous peine de damnation, lorsque la chose est de conséquence. S’il est difficile au riche d’entrer dans le ciel avec son propre bien, comment y entrera-t-il avec le bien d’autrui? Restituez tout, restituez même les dommages; restituez du moins tout ce que vous pourrez; restituez d’abord et ne différez pas. Si vous différez, vous ne restituerez jamais. Restituez à qui vous devez, et demandez avis pour exécuter vos restitutions. Si vous ne le pouvez aujourd’hui, épargnez, ménagez, retranchez vos dépenses, vendez quelques fonds. Il vaut mieux incommoder votre famille, perdre même vos biens, s’il le faut, que de perdre votre ame. Si absolument vous ne pouvez restituer, ni le tout ni une partie, priez du moins pour ceux à qui vous avez fait tort.
 
Le bien usurpé ne profite point. Cent sous mal acquis en feront perdre mille. Si vous laissez du bien d’autrui dans votre famille, il fera périr vos biens, fera damner vos enfans, et vous avec eux. On pèche bien plus griévement, lorsque l’on prend, sachant que l’on ne pourra pas rendre.
 
\secundo Quant à vous, si l’on vous a fait du tort, souffrez avec patience; ne vous vengez point. Vous pouvez vous faire rendre justice, mais que ce soit sans rancune et sans chicane. Il vaut mieux, après tout, qu’on vous fasse tort que de faire tort aux autres. L’injustice qu’on vous fait ne vous rend pas coupable, elle est pour vous un sujet de mérite, si vous avez assez de patience pour souffrir, et assez de charité pour pardonner.
 
Celui qui vous fait tort se fait à lui-même plus de tort qu’à vous, puisqu’il s’expose à perdre son âme. Ne souhaitez jamais sur sa conscience ou sur sa damnation le tort qu’il vous fait: il se damnera assez sans vous. Il est plus digne de votre compassion que de votre colère; vous devez avoir plus de regret de la perte de son âme que de la perte de vos biens.
 
La compensation est ici le plus dangereux piège et le plus ordinaire; la cupidité y a ordinairement plus de part que l’équité et que la raison; elle met souvent la conscience dans d’étranges embarras.
 
\section{Résolutions}
 
\primo J’examinerai sérieusement si je ne me suis point rendu coupable d’injustice envers le prochain, soit en négligeant les devoirs de mon état, soit par fraude, subtilité, mensonge, soit par des acquisitions illégitimes, etc.

\secundo Et si ma conscience me reproche quelque chose de ce genre, je me hâterai de réparer mes torts.

\tertio S’il s’élève dans mon cœur quelque doute en cette matière, je prendrai de suite des mesures pour éclaircir ces doutes, et mettre ordre à ma peine. 

Mon Dieu, ne permettez pas que je m’aveugle sur mes injustices passées; faites-moi la grâce de les réparer comme je le dois, et de n’en plus commettre à l’avenir.
 
\end{document}
